\section{Conclusion}
\label{sec:conclusion}

Overall coverage hides unreachable classes. A narrowing rule can contain the gold label for almost every test instance and still never offer a rare class. Rule: inspect per-class coverage, not overall coverage, before trusting a narrowing rule.

At matched set size the narrowing rule matters under a skewed labelled pool and not on balanced data, and under a 100-fold skew the wrong rule costs several times what narrowing gains over few-shot prompting. Rule: calibrate per class whenever the labelled pool is skewed.

The LLM almost never recovers a label that is missing from the candidate set, so the set is a hard constraint. Rule: choose $\alpha$ with the rare classes in mind, since their coverage bounds macro-F1.

Narrowing pays off modestly on short texts with informative labels and more when label names carry no meaning, and for the encoder recipes we tested, fine-tuning wins on a plentiful balanced pool. Rule: prompt with a class-conditionally narrowed list when labels are skewed or scarce, and fine-tune when they are plentiful and balanced.
